\begin{ouverture}[courbe]{1}{Nombres rationnels et fractions}{Partie 1 · Nombres et calculs}
{Vocabulaire · Priorités · Fractions irréductibles · Problèmes}
\objectif{Calculer en respectant les priorités et les signes.}
\objectif{Rendre une fraction irréductible et calculer exactement.}
\objectif{Résoudre un problème faisant appel à des fractions.}
\prerequis{Tables de multiplication, nombres relatifs et opérations sur les fractions.}
\end{ouverture}
\pagestyle{otMain}\essentiel\fontsize{11}{15}\selectfont
\essSection{1}{Vocabulaire et priorités}
\begin{definition}Un quotient est le résultat d’une division, ou l’écriture de cette division. Une fraction est le quotient $a/b$ de deux entiers, avec $b\ne0$ : elle désigne à la fois un nombre et une écriture. Un nombre rationnel est un nombre égal au quotient de deux entiers.\end{definition}
\begin{propriete}Les parenthèses sont prioritaires, puis les puissances, les produits et quotients, enfin les sommes et différences. À priorité égale, on calcule de gauche à droite.\end{propriete}
\begin{exemple}$-4+3\times(5-7)=-4+3\times(-2)=-10$. Pour $(-3)^2$, le carré porte sur le nombre $-3$ ; dans $-3^2$, le signe moins est extérieur au carré.\end{exemple}
\begin{savoirfaire}[Gérer les signes]Un produit ou un quotient de deux nombres de même signe est positif ; de signes contraires, il est négatif. Pour une somme, ces règles ne s’appliquent pas : comparer les valeurs et les signes.\end{savoirfaire}
\begin{definition}[Ensembles de nombres]$\mathbb N$ désigne les entiers naturels, $\mathbb Z$ les entiers relatifs, $\mathbb D$ les nombres décimaux et $\mathbb Q$ les nombres rationnels. Chacun est contenu dans le suivant. Ainsi $-3\in\mathbb Z$, $2{,}5\in\mathbb D$ et $1/3\in\mathbb Q$, mais $1/3$ n’est pas décimal.\end{definition}
\begin{attention}On ne divise jamais par zéro. Une égalité doit relier des nombres égaux : ne pas enchaîner un résultat partiel avec tout le calcul initial.\end{attention}
\clearpage
\essSection{2}{Calculer avec des fractions}
\begin{propriete}Pour additionner, utiliser un dénominateur commun. Pour multiplier, multiplier les numérateurs et les dénominateurs. Pour diviser par une fraction non nulle, multiplier par son inverse.\end{propriete}
\begin{exemple}
\[\frac34-\frac56=\frac9{12}-\frac{10}{12}=-\frac1{12},
\qquad \frac23\div\frac57=\frac23\times\frac75=\frac{14}{15}.\]
\end{exemple}
\begin{definition}Une fraction est irréductible lorsque son numérateur et son dénominateur n’ont pas d’autre diviseur commun positif que 1.\end{definition}
\begin{savoirfaire}[Rendre irréductible]Simplifier successivement par des diviseurs communs, en s’aidant des critères de divisibilité, ou décomposer les deux termes en facteurs premiers. $\dfrac{84}{126}=\dfrac{42}{63}=\dfrac{14}{21}=\dfrac23$ : on a simplifié par 2, par 3, puis par 7.\end{savoirfaire}
\begin{savoirfaire}[Une fraction d’une quantité]Prendre $3/5$ de 40 revient à calculer $40\times3/5=24$. « Les $2/3$ du reste » impose de calculer d’abord le reste : le tout de référence a changé. Un schéma en barres aide à suivre ces parts.\end{savoirfaire}
\begin{attention}On simplifie des facteurs dans un produit, jamais des termes isolés d’une somme. Ainsi $(6+3)/3=3$, et non $6+1$.\end{attention}
\clearpage\rubrique{Découverte}
\begin{blocLivre}{ex-1-1}\begin{exercice}[1]{Une température}
Le matin, il fait $-6\,^\circ$C. La température augmente de 11 degrés puis baisse de 4 degrés. Quelle est la température finale ?
\end{exercice}\end{blocLivre}
\begin{blocLivre}{ex-1-2}\begin{exercice}[1]{Respecter les priorités}
Calculer $A=18-3(2+5)$ et $B=(18-3)(2+5)$.
\end{exercice}\end{blocLivre}
\begin{blocLivre}{ex-1-3}\begin{exercice}[1]{Deux écritures, un nombre}
Compléter $\dfrac5{8}=\dfrac{\,?\,}{40}$ et $\dfrac{42}{63}=\dfrac2{\,?\,}$.
\end{exercice}\end{blocLivre}
\begin{blocLivre}{ex-1-4}\begin{exercice}[1]{Une somme exacte}
Calculer $\dfrac7{12}+\dfrac5{18}$ et donner une fraction irréductible.
\end{exercice}\end{blocLivre}
\begin{blocLivre}{ex-1-5}\begin{exercice}[1]{Produit et quotient}
Calculer $\dfrac{-9}{14}\times\dfrac7{15}$ puis $\dfrac35\div\dfrac{-9}{10}$.
\end{exercice}\end{blocLivre}

\clearpage\rubrique{Entraînement}
\begin{blocLivre}{ex-1-6}\begin{exercice}[2]{Simplifications successives}
Rendre irréductibles $\dfrac{150}{225}$ et $\dfrac{126}{198}$ par simplifications successives. Les deux fractions sont-elles égales ?
\end{exercice}\end{blocLivre}
\begin{blocLivre}{ex-1-7}\begin{exercice}[2]{Comparer sans arrondir}
Ranger $-\dfrac34$, $-\dfrac56$ et $-\dfrac7{12}$ dans l’ordre croissant.
\end{exercice}\end{blocLivre}
\begin{blocLivre}{ex-1-8}\begin{exercice}[2]{Une expression à étages}
Calculer exactement $C=\dfrac{\frac23+\frac16}{\frac54-\frac12}$.
\end{exercice}\end{blocLivre}
\begin{blocLivre}{ex-1-9}\begin{exercice}[2]{Des ensembles de nombres}
Pour chacun des nombres $-7$ ; $0$ ; $2{,}5$ ; $\dfrac{12}{4}$ ; $-\dfrac53$, donner le plus petit des ensembles $\mathbb N$, $\mathbb Z$, $\mathbb D$, $\mathbb Q$ qui le contient.
\end{exercice}\end{blocLivre}
\begin{blocLivre}{ex-1-10}\begin{exercice}[2]{Une réduction trompeuse}
Un élève écrit $(8+4)/4=8+1$. Expliquer son erreur et calculer correctement.
\end{exercice}\end{blocLivre}

\clearpage\rubrique{Problèmes}
\begin{blocLivre}{ex-1-11}\begin{exercice}[3]{Deux étapes de lecture}
Un lecteur lit $2/5$ d’un livre le lundi, puis $1/3$ de ce qui reste le mardi. Quelle fraction du livre reste-t-il ? Combien de pages pour un livre de 270 pages ?
\end{exercice}\end{blocLivre}
\begin{blocLivre}{ex-1-12}\begin{exercice}[3]{Des bouquets vendus}
Une fleuriste vend les $2/3$ de ses bouquets le matin, puis les $3/4$ du reste l’après-midi. Le soir, il lui reste 9 bouquets. Combien en avait-elle au début ?
\end{exercice}\end{blocLivre}
\begin{blocLivre}{ex-1-13}\begin{exercice}[3]{Une réserve d’eau}
Une cuve est remplie aux $7/8$. On consomme $2/7$ de l’eau présente, puis on ajoute 30 L. La cuve est alors remplie aux $3/4$. Quelle est sa capacité ?
\end{exercice}\end{blocLivre}
\begin{blocLivre}{ex-1-14}\begin{exercice}[3]{Des votes et des absents}
Dans une classe, $5/6$ des élèves sont présents. Parmi eux, $3/5$ choisissent le projet A. Le projet A reçoit 15 voix. Combien y a-t-il d’élèves et d’absents ?
\end{exercice}\end{blocLivre}
\begin{blocLivre}{ex-1-15}\begin{exercice}[3]{Trouver toutes les possibilités}
On cherche les entiers naturels $n$ tels que $n/12$ soit strictement supérieur à $2/3$ et strictement inférieur à $5/4$. Les donner tous.
\end{exercice}\end{blocLivre}

\clearpage
